% =====================================================================
\chapter{Regresión lineal y mínimos cuadrados}
\label{tema:2}
% =====================================================================

Este tema estudia el modelo más sencillo de aprendizaje supervisado, la regresión lineal ajustada por mínimos cuadrados. Se deducen las ecuaciones normales, se interpretan como una proyección ortogonal y se justifica la pérdida cuadrática como máxima verosimilitud con ruido gaussiano. La exposición sigue el capítulo~1 de los apuntes de CS229~\cite{cs229}, las secciones~4.2 y~4.3 de Strang~\cite{strang} y el curso 18.06SC del MIT~\cite{mit1806}, la sección~3.1 de Bishop~\cite{bishop} y la sección~3.1 de Zhang \emph{et al.}~\cite{d2l}.

\section{El modelo lineal y la matriz de diseño}

\begin{definicion}[Regresión lineal]
Un \emph{modelo de regresión lineal} predice $\hat y=w^{\top}x+b$, con \emph{pesos} $w\in\R^{d}$ y \emph{término independiente} $b\in\R$. Añadiendo a cada entrada una componente constante, $\tilde x=(1,x)\in\R^{d+1}$, y escribiendo $\theta=(b,w)$, el modelo queda $\hat y=\theta^{\top}\tilde x$. En lo sucesivo se omite la tilde.
\end{definicion}

\begin{definicion}[Matriz de diseño y suma de cuadrados]
Dados $n$ ejemplos, la \emph{matriz de diseño} $X\in\R^{n\times(d+1)}$ tiene por fila $i$ el vector $x_i^{\top}$ (con su primera componente igual a~1), y $y=(y_1,\dots,y_n)^{\top}$. El vector de predicciones es $\hat y=X\theta$, y la \emph{suma de cuadrados de los residuos} es
\[
J(\theta)=\norm{y-X\theta}^{2}=\sum_{i=1}^{n}\bigl(y_i-\theta^{\top}x_i\bigr)^{2}.
\]
El estimador de \emph{mínimos cuadrados} es cualquier $\hat\theta$ que minimice $J$; se escribe $\SSE=J(\hat\theta)$.
\end{definicion}

\begin{observacion}[Funciones base]
El modelo es lineal en $\theta$, no necesariamente en $x$. Si se sustituye $x$ por un vector de \emph{funciones base} $\phi(x)$, por ejemplo $\phi(x)=(1,x,x^{2},\dots,x^{p})$, todo lo que sigue se aplica sin cambios con $X_{ij}=\phi_j(x_i)$~\cite[§3.1]{bishop}. Así se ajustan los polinomios del \cref{tema:5}.
\end{observacion}

\section{Las ecuaciones normales}

\begin{lema}[Gradientes de formas lineales y cuadráticas]\label{lem:gradientes}
Sean $a\in\R^{m}$ y $A\in\R^{m\times m}$. Entonces $\nabla_\theta(a^{\top}\theta)=a$ y $\nabla_\theta(\theta^{\top}A\theta)=(A+A^{\top})\theta$, que vale $2A\theta$ si $A$ es simétrica.
\end{lema}

\begin{proof}
$\partial(\sum_j a_j\theta_j)/\partial\theta_k=a_k$. Para la forma cuadrática, $\partial(\sum_{i,j}\theta_iA_{ij}\theta_j)/\partial\theta_k=\sum_j A_{kj}\theta_j+\sum_i\theta_iA_{ik}=(A\theta)_k+(A^{\top}\theta)_k$ (módulo~4).
\end{proof}

\begin{teorema}[Ecuaciones normales]\label{teo:normales}
Un vector $\hat\theta$ minimiza $J(\theta)=\norm{y-X\theta}^{2}$ si y solo si cumple las \emph{ecuaciones normales}
\[
X^{\top}X\,\hat\theta=X^{\top}y.
\]
Además, para cualquier solución $\hat\theta$ y cualquier $\theta$, se verifica $J(\theta)=J(\hat\theta)+\norm{X(\theta-\hat\theta)}^{2}$.
\end{teorema}

\begin{proof}
Desarrollando, $J(\theta)=y^{\top}y-2\theta^{\top}X^{\top}y+\theta^{\top}X^{\top}X\theta$. Por el \cref{lem:gradientes}, con $A=X^{\top}X$ simétrica, $\nabla J(\theta)=-2X^{\top}y+2X^{\top}X\theta$, que se anula exactamente en las soluciones de las ecuaciones normales. Para ver que son mínimos globales, sea $\hat\theta$ una solución y $r=y-X\hat\theta$, que cumple $X^{\top}r=0$. Para cualquier $\theta$,
\[
y-X\theta=r-X(\theta-\hat\theta),\qquad
J(\theta)=\norm{r}^{2}-2(\theta-\hat\theta)^{\top}X^{\top}r+\norm{X(\theta-\hat\theta)}^{2}=J(\hat\theta)+\norm{X(\theta-\hat\theta)}^{2}.
\]
Luego $J(\theta)\ge J(\hat\theta)$. Recíprocamente, si $\theta$ es un minimizador, $\nabla J(\theta)=0$ y $\theta$ cumple las ecuaciones normales.
\end{proof}

\begin{lema}\label{lem:rango}
$X^{\top}X$ es simétrica y semidefinida positiva. Es definida positiva, y por tanto invertible, si y solo si las columnas de $X$ son linealmente independientes, es decir, si $\rg(X)=d+1$.
\end{lema}

\begin{proof}
$(X^{\top}X)^{\top}=X^{\top}X$ y $v^{\top}X^{\top}Xv=\norm{Xv}^{2}\ge 0$. Esta cantidad es nula si y solo si $Xv=0$, y existe $v\ne 0$ con $Xv=0$ si y solo si las columnas son dependientes.
\end{proof}

\begin{corolario}[Unicidad]\label{cor:unicidad}
Si $\rg(X)=d+1$, el estimador de mínimos cuadrados es único y vale $\hat\theta=(X^{\top}X)^{-1}X^{\top}y$. Si $\rg(X)<d+1$, hay infinitos minimizadores, todos con las mismas predicciones $X\hat\theta$, y el de norma mínima es $X^{+}y$, donde $X^{+}$ es la pseudoinversa (módulo~3).
\end{corolario}

\begin{proof}
El primer caso es el \cref{lem:rango}. En el segundo, por el \cref{teo:normales}, $\theta$ y $\hat\theta$ son ambos minimizadores si y solo si $\norm{X(\theta-\hat\theta)}^{2}=0$, es decir, si $\theta-\hat\theta\in\ker X$, que es un subespacio no trivial; las predicciones coinciden. La caracterización de la solución de norma mínima mediante la descomposición en valores singulares se demostró en el módulo~3.
\end{proof}

\begin{ejemplo}[Ajuste de una recta]\label{ej:recta}
Sean $x=(0,1,2,3)$ e $y=(1,2,4,5)$. La matriz de diseño tiene filas $(1,x_i)$, y
\[
X^{\top}X=\begin{pmatrix} n & \sum x_i\\ \sum x_i & \sum x_i^{2}\end{pmatrix}=\begin{pmatrix}4&6\\6&14\end{pmatrix},\qquad
X^{\top}y=\begin{pmatrix}\sum y_i\\ \sum x_iy_i\end{pmatrix}=\begin{pmatrix}12\\25\end{pmatrix}.
\]
El determinante es $4\cdot 14-6^{2}=20$ y, por la regla de Cramer, $b=(12\cdot 14-6\cdot 25)/20=0{,}9$ y $w=(4\cdot 25-6\cdot 12)/20=1{,}4$. Las predicciones son $(0{,}9;\,2{,}3;\,3{,}7;\,5{,}1)$, los residuos $r=(0{,}1;\,-0{,}3;\,0{,}3;\,-0{,}1)$ y $\SSE=0{,}2$. Se comprueba que $\sum_i r_i=0$ y $\sum_i x_ir_i=-0{,}3+0{,}6-0{,}3=0$, es decir, $X^{\top}r=0$.
\end{ejemplo}

\begin{observacion}[Cálculo numérico]
En la práctica no se calcula $(X^{\top}X)^{-1}$: se resuelve el sistema con una factorización de Cholesky, o se aplica a $X$ una factorización QR o la descomposición en valores singulares, que son más estables porque el número de condición de $X^{\top}X$ es el cuadrado del de $X$. Con millones de parámetros, como en una red neuronal, no existe fórmula cerrada y se usa el descenso de gradiente (\cref{sec:lms}).
\end{observacion}

\section{Interpretación geométrica}

Las ecuaciones normales equivalen a $X^{\top}(y-X\hat\theta)=0$: el residuo es ortogonal a todas las columnas de $X$. Por tanto, $\hat y=X\hat\theta$ es el punto del espacio de columnas $\mathcal{C}(X)$ más próximo a $y$, su \emph{proyección ortogonal}~\cite[§4.2--4.3]{strang}.

\begin{proposicion}[Matriz de proyección]\label{prop:proyeccion}
Supóngase que $\rg(X)=d+1$ y sea $P=X(X^{\top}X)^{-1}X^{\top}$, de tamaño $n\times n$. Entonces:
\begin{enumerate}
  \item $\hat y=Py$ y $r=(I-P)y$;
  \item $P^{\top}=P$ y $P^{2}=P$;
  \item $\hat y^{\top}r=0$ y, en consecuencia, $\norm{y}^{2}=\norm{\hat y}^{2}+\norm{r}^{2}$;
  \item si $X$ contiene la columna de unos, $\sum_i r_i=0$, y la recta (o el hiperplano) ajustada pasa por el punto medio $(\bar x,\bar y)$.
\end{enumerate}
\end{proposicion}

\begin{proof}
(i) $\hat y=X\hat\theta=X(X^{\top}X)^{-1}X^{\top}y$. (ii) La transpuesta de $(X^{\top}X)^{-1}$ es ella misma por ser $X^{\top}X$ simétrica, luego $P^{\top}=P$; y $P^{2}=X(X^{\top}X)^{-1}(X^{\top}X)(X^{\top}X)^{-1}X^{\top}=P$. (iii) $\hat y^{\top}r=y^{\top}P(I-P)y=y^{\top}(P-P^{2})y=0$; el teorema de Pitágoras da la igualdad de normas. (iv) La primera ecuación normal es el producto de la columna de unos por $r$: $\sum_i r_i=0$. Dividiendo por $n$, $\bar y=\hat\theta^{\top}\bar x$, donde $\bar x$ es la media de las filas de $X$.
\end{proof}

\begin{figure}[h]
\centering
\begin{tikzpicture}[font=\small,>={Stealth[length=2.2mm]}]
  \fill[gray!12] (-0.6,-0.9) -- (5.0,-0.9) -- (6.4,1.2) -- (0.8,1.2) -- cycle;
  \node at (5.3,-0.6) {$\mathcal{C}(X)$};
  \coordinate (O) at (0,0);
  \coordinate (P) at (3.6,0.5);
  \coordinate (Y) at (3.6,3.0);
  \draw[->,thick] (O) -- (Y) node[above left] {$y$};
  \draw[->,thick] (O) -- (P) node[below right] {$\hat y=X\hat\theta$};
  \draw[->,thick,dashed] (P) -- (Y) node[midway,right] {$r=y-\hat y$};
  \draw (3.6,0.8) -- (3.32,0.76) -- (3.32,0.46);
  \draw[->] (O) -- (1.6,-0.6) node[below] {$X_{:,1}$};
  \draw[->] (O) -- (1.3,0.9) node[above] {$X_{:,2}$};
\end{tikzpicture}
\caption{Mínimos cuadrados como proyección ortogonal: $\hat y$ es el punto del espacio de columnas más próximo a $y$, y el residuo $r$ es perpendicular a ese espacio.}
\label{fig:proyeccion}
\end{figure}

En el \cref{ej:recta}, $\norm{y}^{2}=46$, $\norm{\hat y}^{2}=45{,}8$ y $\norm{r}^{2}=0{,}2$, de acuerdo con la \cref{prop:proyeccion}(iii).

\section{Descenso de gradiente: el algoritmo LMS}
\label{sec:lms}

Cuando $n$ o $d$ son grandes, se minimiza $J$ iterativamente. Con la pérdida media $\frac{1}{2n}\norm{y-X\theta}^{2}$, el gradiente es $-\frac1n X^{\top}(y-X\theta)$, y el descenso de gradiente con tasa $\alpha>0$ actualiza
\[
\theta\leftarrow\theta+\frac{\alpha}{n}\sum_{i=1}^{n}\bigl(y_i-\theta^{\top}x_i\bigr)x_i.
\]
Si en cada paso se usa un solo ejemplo, se obtiene la \emph{regla LMS} (\emph{least mean squares}) o de Widrow-Hoff, $\theta\leftarrow\theta+\alpha\,(y_i-\theta^{\top}x_i)\,x_i$: la corrección es proporcional al error cometido en ese ejemplo~\cite[§1.1]{cs229}. Como $J$ es convexa (su hessiano $2X^{\top}X$ es semidefinido positivo, \cref{lem:rango}), el descenso de gradiente converge al mínimo global si la tasa es suficientemente pequeña (módulo~7).

\section{Interpretación probabilística}

La pérdida cuadrática no es una elección arbitraria: es la que resulta de suponer ruido gaussiano~\cite[§1.3]{cs229}, \cite[§3.1.1]{bishop}.

\begin{definicion}[Modelo lineal gaussiano]
Las respuestas cumplen $y_i=\theta^{\top}x_i+\varepsilon_i$, donde los errores $\varepsilon_1,\dots,\varepsilon_n$ son independientes con distribución $\N(0,\sigma^{2})$. Equivalentemente, $y_i\mid x_i\sim\N(\theta^{\top}x_i,\sigma^{2})$.
\end{definicion}

\begin{teorema}[Mínimos cuadrados y máxima verosimilitud]\label{teo:mv-gauss}
En el modelo lineal gaussiano, la log-verosimilitud negativa es
\[
-\ln L(\theta,\sigma^{2})=\frac{n}{2}\ln(2\pi\sigma^{2})+\frac{1}{2\sigma^{2}}\norm{y-X\theta}^{2}.
\]
Para cualquier $\sigma^{2}>0$, el estimador de máxima verosimilitud de $\theta$ es el de mínimos cuadrados, y el de $\sigma^{2}$ es $\hat\sigma^{2}=\SSE/n$.
\end{teorema}

\begin{proof}
La densidad de $y_i$ es $(2\pi\sigma^{2})^{-1/2}\exp\bigl(-(y_i-\theta^{\top}x_i)^{2}/(2\sigma^{2})\bigr)$; por independencia, la verosimilitud es el producto y su logaritmo la suma, lo que da la expresión del enunciado. El primer sumando no depende de $\theta$ y el segundo es $J(\theta)$ multiplicado por la constante positiva $1/(2\sigma^{2})$, luego ambos se minimizan en el mismo $\hat\theta$. Sustituyendo $\hat\theta$ y derivando respecto a $s=\sigma^{2}$: $\frac{n}{2s}-\frac{\SSE}{2s^{2}}=0$, es decir, $s=\SSE/n$; la derivada pasa de negativa a positiva en ese punto, que es por tanto un mínimo.
\end{proof}

En el \cref{ej:recta}, $\hat\sigma^{2}=0{,}2/4=0{,}05$. El razonamiento se generaliza: cada supuesto sobre la distribución de $y$ produce su pérdida de máxima verosimilitud (\cref{tab:mv}). Con ruido de Laplace, por ejemplo, la pérdida es la suma de valores absolutos (\cref{eje:2-laplace}).

\begin{table}[h]
\centering\small
\caption{Pérdida de máxima verosimilitud según la distribución supuesta para $y$.}\label{tab:mv}
\begin{tabular}{@{}llll@{}}
\toprule
Respuesta & Distribución & Pérdida & Modelo \\
\midrule
real & normal & cuadrática & regresión lineal (\cref{tema:2}) \\
binaria & Bernoulli & entropía cruzada binaria & regresión logística (\cref{tema:3}) \\
una de $K$ clases & categórica & entropía cruzada & regresión softmax (\cref{tema:4}) \\
token siguiente & categórica sobre $\V$ & entropía cruzada & capa de salida de un LLM \\
\bottomrule
\end{tabular}
\end{table}

\begin{observacion}[Física]
$\frac12\sum_i(y_i-\hat y_i)^{2}$ es la energía elástica de $n$ muelles verticales de constante unidad que unen cada dato con la recta. La recta de mínimos cuadrados es la posición de equilibrio, y las ecuaciones normales expresan que la fuerza total y el momento total de los muelles son nulos.
\end{observacion}

\begin{lecturas}
\item Ng y Ma~\cite{cs229}, cap.~1: §1.1 «LMS algorithm», §1.2 «The normal equations» y §1.3 «Probabilistic interpretation».
\item Strang~\cite{strang}, §4.2 «Projections» y §4.3 «Least Squares Approximations»; MIT 18.06SC~\cite{mit1806}, sesión~2.3.
\item Bishop~\cite{bishop}, §3.1.1 «Maximum likelihood and least squares» y §3.1.2 «Geometry of least squares».
\item MIT 6.390~\cite{mit6390}, cap.~2, §2.5 «Analytical solution: ordinary least squares».
\item Zhang \emph{et al.}~\cite{d2l}, §3.1 «Linear Regression», en particular §3.1.3.
\end{lecturas}

\section*{Ejercicios del tema 2}
\addcontentsline{toc}{section}{Ejercicios del tema 2}

\begin{ejercicio}{1}{eje:2-rango}
Con los datos $x=(0,1,2)$, $y=(0,1,3)$ se construye por error una matriz de diseño con tres columnas: unos, $x$ y otra vez $x$. (a) ¿Es invertible $X^{\top}X$? (b) Describa el conjunto de minimizadores sabiendo que, con una sola columna $x$, la solución es $b=-1/6$, $w=3/2$. (c) Proponga dos formas de resolver el problema y halle la solución de norma mínima.
\fuente{adaptado de Zhang \emph{et al.}~\cite{d2l}, §3.1, ejercicio~4, apartados (a) y (b).}
\end{ejercicio}

\begin{ejercicio}{2}{eje:2-recta}
Ajuste por mínimos cuadrados una recta a $x=(0,1,2)$, $y=(0,1,3)$. Calcule $X^{\top}X$, $X^{\top}y$, $\hat\theta$, los residuos, el $\SSE$ y el estimador de máxima verosimilitud $\hat\sigma^{2}$.
\end{ejercicio}

\begin{ejercicio}{2}{eje:2-strang1}
Se mide $b$ en los instantes $t=-1,1,2$ y se obtiene $b=(1,3,11)$. (a) Ajuste la recta $b=C+Dt$ por mínimos cuadrados. (b) Calcule las alturas $p=A\hat x$ y el error $e=b-p$. (c) Compruebe que $e$ es ortogonal a las dos columnas de $A$. (d) ¿Qué recta se obtiene si los datos son exactamente $e$?
\fuente{adaptado de MIT 18.06SC~\cite{mit1806}, sesión~2.3 «Projection Matrices and Least Squares», problemas 16.1--16.3 (Strang~\cite{strang}, §4.3).}
\end{ejercicio}

\begin{ejercicio}{2}{eje:2-strang2}
Ajuste la mejor recta $C+Dt$ a $b=(7,4,2,1,1)$ en los instantes $t=-2,-1,0,1,2$. Explique por qué $A^{\top}A$ resulta diagonal y calcule el $\SSE$.
\fuente{adaptado de MIT 18.06SC~\cite{mit1806}, sesión~2.3, problema~16.6 (Strang~\cite{strang}, §4.3).}
\end{ejercicio}

\begin{ejercicio}{2}{eje:2-laplace}
Supóngase que $y_i=b+\varepsilon_i$ con errores independientes de Laplace, de densidad $p(\varepsilon)=\frac{1}{2s}e^{-\abs{\varepsilon}/s}$ y escala $s>0$. (a) Escriba la log-verosimilitud negativa de $(b,s)$. (b) Halle los estimadores de máxima verosimilitud para los datos $y=(1,2,4,9,14)$. (c) Explique por qué, en el modelo lineal $y_i=\theta^{\top}x_i+\varepsilon_i$ con este ruido, no existe una fórmula cerrada como las ecuaciones normales.
\fuente{adaptado de Zhang \emph{et al.}~\cite{d2l}, §3.1, ejercicio~5, apartados (a) y (b).}
\end{ejercicio}

\begin{ejercicio}{3}{eje:2-traza}
Sea $X\in\R^{n\times m}$ de rango $m$ y $P=X(X^{\top}X)^{-1}X^{\top}$. (a) Demuestre que $I-P$ es también simétrica e idempotente y que $(I-P)X=0$. (b) Demuestre que $\tr P=m$, el número de parámetros. (c) Demuestre que los autovalores de $P$ son 0 o~1. (d) Para la matriz de diseño del \cref{ej:recta}, compruebe que la diagonal de $P$ es $(0{,}7;\,0{,}3;\,0{,}3;\,0{,}7)$.
\fuente{adaptado de Strang~\cite{strang}, §4.2 «Projections» (propiedades $P^{\top}=P$ y $P^{2}=P$).}
\end{ejercicio}
